
% --------------------------------------------------
\section{Өгөгдөл}
\label{sec:data}
% --------------------------------------------------

Энэхүү судалгаанд ковариацын матрицын агшилтын
үнэлэгчдийн гүйцэтгэлийг хөрөнгүүдийн хоорондын
корреляцын ялгаатай бүтэцтэй багцуудад харьцуулах
зорилгоор АНУ-ын S\&P 500 индексийн түүхэн
бүрэлдэхүүнээс сонгосон хувьцааны сарын өгөөжийг
ашигласан.


Нэр дэвшигч хувьцаануудыг S\&P 500 индексийн
түүхэн бүрэлдэхүүний нээлттэй өгөгдлийн сангаас
2012 оны 5 дугаар сарын 30-ны өдрийн байдлаар
сонгосон \citep{chinobing_sp500}.
Тухайн жагсаалтад 453 хувьцааны код
байсан бөгөөд 2012 оны 5 дугаар сараас 2026 оны
9 дүгээр сар хүртэлх сарын эцсийн үнийн мэдээлэл
бүрэн байгаа 337 хувьцааг шинжилгээнд хамруулсан.

Хувьцааны үнийн өгөгдлийг Yahoo Finance мэдээллийн
сангаас Python програмчлалын хэлний
\texttt{yfinance} сангаар татаж,
\texttt{auto\_adjust=True} тохиргоогоор
ногдол ашиг болон хувьцаа хуваалтын нөлөөг
харгалзан үзсэн. Хувьцаа тус бүрийн сарын
энгийн өгөөжийг
\[
R_{i,t}
=
\frac{P_{i,t}}{P_{i,t-1}}-1
\]
гэж тооцсон. Энд $P_{i,t}$ нь $i$-р хувьцааны
$t$ сарын эцсийн тохируулгатай хаалтын үнэ юм.

Судалгаанд 2012 оны 6 дугаар сараас 2026 оны
9 дүгээр сар хүртэлх нийт $n=172$ сарын өгөөжийг
ашигласан. Багцуудыг сонгохдоо 2012 оны
6 дугаар сараас 2022 оны 5 дугаар сар хүртэлх
эхний 120 сарын өгөөжид үндэслэн хөрөнгүүдийн
хоорондын корреляцыг тооцсон.

Үндсэн шинжилгээнд дундаж хос корреляц өндөр,
дунд болон бага байх тус бүр 10 хувьцаатай,
хоорондоо давхцахгүй гурван багцыг Monte Carlo
хайлтын аргаар байгуулсан. Дунд корреляцын
зорилтот түвшнийг 337 хувьцааны бүх хосын
дундаж корреляцаар тодорхойлсон.

\begin{table}[htbp]
	\centering
	\caption{Үндсэн шинжилгээнд сонгосон
		багцуудын шинж чанар}
	\label{tab:portfolio_selection}
	\small
	\begin{tabular}{lccc}
		\hline
		Үзүүлэлт
		& Өндөр
		& Дунд
		& Бага\\
		\hline
		Хөрөнгийн тоо
		& 10 & 10 & 10\\
		Дундаж корреляц
		& 0.5570 & 0.3036 & 0.1015\\
		Нөхцөлт тоо
		& 72.99 & 16.82 & 27.45\\
		\hline
	\end{tabular}
\end{table}

Өндөр корреляцтай багцад CCL, CME, COF,
CTAS, FLS, MET, PRU, SWK, UNM, WFC;
дунд корреляцтай багцад AMGN, AMT, BEN,
DTE, EXPD, MO, MRK, PG, SNA, UPS;
бага корреляцтай багцад BAX, BIIB, CTSH,
GNW, HRB, MTB, NEM, QCOM, RRC, WEC
хувьцаанууд сонгогдсон.

Үр дүн багцын бүрэлдэхүүнээс хамаарах эсэхийг
шалгах зорилгоор нэмэлт 150 багц сонгосон.
Үүнд өндөр, дунд, бага корреляцын ангилал
тус бүрээс 50 багц хамруулсан. Багцуудыг
эхний 120 сарын өгөөжид үндэслэн Monte Carlo
хайлтын аргаар сонгож, нэг ангиллын хоёр
багцад хамгийн ихдээ зургаан ижил хувьцаа
байх нөхцөл тавьсан. Сонголтын дэлгэрэнгүй
аргачлалыг
\ref{subsubsec:portfolio_robustness}
хэсэгт тайлбарласан.

Үндсэн шинжилгээнд ковариацын матрицыг
$T=120$ сарын шилжих цонхоор үнэлж,
2022 оны 6 дугаар сараас 2026 оны
9 дүгээр сар хүртэлх $M=52$ сарын
түүврээс гадуурх (OOS) гүйцэтгэлийг
харьцуулсан. Сонгосон багц бүрийн хувьцааны
бүрэлдэхүүнийг OOS хугацаанд тогтмол хадгалж,
багцын оновчтой жинг сар бүр шинэчилсэн.

Судалгааны нэг хязгаарлалт нь шинжилгээний
бүх хугацаанд үнийн мэдээлэл бүрэн байгаа
хувьцаануудыг сонгосон явдал юм. Энэ нь
survivorship bias болон ирээдүйн мэдээлэлд
тулгуурласан сонголтын хазайлт үүсгэж
болзошгүй тул үр дүнг тайлбарлахдаа
уг хязгаарлалтыг харгалзан үзнэ.
